% !TeX root = ../main.tex
% =================================================================================
\chapter{Related Work}\label{chap:related-work}
% =================================================================================

In this chapter, we review work related to this thesis.
First, \cref{sec:related-work:method} describes our literature search methodology.
Next, \cref{sec:related-work:first,sec:related-work:second} review related work in the areas of \ldots and \ldots
Last, \cref{sec:related-work:discussion} compares the related work and delimits our work from it.

% ---------------------------------------------------------------------------------
\section{Literature Search Methodology}\label{sec:related-work:method}
% ---------------------------------------------------------------------------------

Describe how you identified related work so that others can reproduce it:
Which search engines and databases did you use (e.g., Google Scholar, ACM Digital Library, IEEE Xplore)?
Which search queries?
How many results did you consider, and which inclusion and exclusion criteria did you apply (title, abstract, full text)?
Did you perform forward and backward snowballing?

% ---------------------------------------------------------------------------------
\section{First Area of Related Work}\label{sec:related-work:first}
% ---------------------------------------------------------------------------------

We identify related work in this area using the following search terms:
\begin{itemize}
  \item \texttt{first search term}
  \item \texttt{second search term}
\end{itemize}

Summarize each relevant work in a few sentences and explain how it relates to your work.
Group related works instead of listing them one after the other.
Explain why they do not (fully) solve the problem of your thesis, e.g., because they make more limiting assumptions, address a different aspect, or lack an evaluation.

% ---------------------------------------------------------------------------------
\section{Second Area of Related Work}\label{sec:related-work:second}
% ---------------------------------------------------------------------------------

Proceed in the same way for every further area of related work.

% ---------------------------------------------------------------------------------
\section{Discussion}\label{sec:related-work:discussion}
% ---------------------------------------------------------------------------------

Compare the related work and delimit your own work from it.
A table that compares the approaches along important criteria, such as \cref{tab:related-work:comparison}, helps the reader.

\begin{table}[htbp]
  \centering
  \caption{Comparison of related approaches (example)}
  \label{tab:related-work:comparison}
  \begin{tabularx}{0.85\textwidth}{@{}Xcccc@{}}
    \toprule
    Approach            & Criterion 1 & Criterion 2 & Criterion 3 & Criterion 4 \\
    \midrule
    Approach A~\cite{Basili1994}   & \cmark & \xmark & \xmark & \cmark \\
    Approach B~\cite{Runeson2009}  & \cmark & \cmark & \xmark & \xmark \\
    This thesis         & \cmark & \cmark & \cmark & \cmark \\
    \bottomrule
  \end{tabularx}
\end{table}
